% ECONOMICS_PROBLEM: {"id": "C32-504", "title": "Economic interpretation of statistically identified shocks", "jel_code": "C32", "source_id": "OP-0604"}
% FIRST_POSTED: 2026-10-09
% ATTEMPT_STATUS: partial
% ENTRY_KIND: research_attempt
\documentclass[11pt,a4paper]{article}
\usepackage[T1]{fontenc}
\usepackage[utf8]{inputenc}
\usepackage[margin=27mm]{geometry}
\usepackage{amsmath,amssymb,amsthm,mathtools}
\usepackage[hidelinks]{hyperref}
\setlength{\parindent}{0pt}
\setlength{\parskip}{5pt}
\newtheorem{theorem}{Theorem}[section]
\newtheorem{proposition}[theorem]{Proposition}
\newtheorem{lemma}[theorem]{Lemma}
\newcommand{\R}{\mathbb R}
\newcommand{\E}{\mathbb E}
\DeclareMathOperator{\Var}{Var}
\DeclareMathOperator{\Cov}{Cov}
\begin{document}
\section*{Statistical shock identification and economic labeling}
\section{A finite labeling problem that can be solved exactly}
Higher-moment identification and economic interpretation answer different questions. This attempt gives an exact algorithm when the statistical ambiguity consists of signed permutations, derives robust proxy-label conditions, and shows why a singleton response need not imply a unique economic label. Continuous rotational ambiguity and uncertain economic exclusions remain substantive limitations.

Assume first that the population innovation matrix is known to belong to
\[
 \mathcal B=\{B_0PD:\ P\text{ a permutation matrix},\quad
 D=\operatorname{diag}(s_1,\ldots,s_m),\ s_j\in\{-1,1\}\},
\]
where $B_0$ is nonsingular. This is an assumption about the full compatible set, not an automatic consequence of observing non-Gaussian residuals. For each candidate impact vector $b=\pm B_0e_j$, define its whole response path
\[
 r(b)=(e_y'A_0b,\ldots,e_y'A_Hb).
\]
Let $\mathcal R$ be specified predicates involving columns, observed external moments, and timing. If restrictions concern other labels as well, test the entire signed-permuted matrix, not just one column.
\begin{proposition}
Enumerating all $2^m m!$ signed permutations and retaining the labeled columns that satisfy $\mathcal R$ produces exactly the finite joint response-path set. Empty output means incompatibility. A unique response path occurs precisely when all retained vectors $r(b)$ coincide.
\end{proposition}
\begin{proof}
Every compatible structural matrix occurs in the enumeration by assumption. Retention is exactly the declared economic restriction. Conversely each retained matrix belongs to $\mathcal B$ and satisfies the restriction, hence is admissible. Taking its path gives both inclusions. The singleton assertion follows directly, with no separate identification claim for the matrix.
\end{proof}
If all restrictions are columnwise and no assignment restriction links different labels, only $2m$ candidates need testing for one label. Global one-to-one labeling constraints require the larger assignment problem. Efficiency is secondary here: the finite algorithm is exact, but may be expensive for many shocks.

The distinction between a unique response and a unique label matters. Take $B_0=I_2$, and let the reported outcome be the sum of two observables, so its selector is $(1,1)'$. At horizon zero both positive columns have response one. If both are economically admissible, the response set is a singleton although the label is ambiguous. If a coordinate selector is required, add a third observable and use a nonsingular matrix with first row $(1,1,0)$; the first two columns still have the same first-coordinate impact. At later horizons the paths can separate. Thus identification of one scalar functional is weaker than identification of the shock's institutional meaning.

\section{An external proxy gives testable labeling restrictions}
Let $c=E[uZ]$ for a scalar external measurement $Z$, and calculate
$d=B_0^{-1}c$. For candidate signed permutation $B=B_0PD$,
\[
 E[\varepsilon Z]=DP'd. \tag{1}
\]
An exact proxy exclusion for target index $j$ requires every nontarget coordinate of (1) to vanish and its target coordinate to be positive. Therefore, before labels are permuted, exactly one nonzero coordinate of $d$ is necessary and sufficient for a unique relevant column direction, with its sign fixed by positive relevance. If $d=0$, relevance fails; if two entries are nonzero, the exact exclusion is incompatible with every one-shock label in this model. This is more informative than assigning the largest correlation and silently discarding the others.

Allow instead $|E[Z\varepsilon_k]|\le\epsilon$ for nontarget shocks and $E[Z\varepsilon_j]\ge\kappa>0$ for the target. A signed original column $sB_0e_j$ is admissible exactly when
\[
 sd_j\ge\kappa,\qquad \max_{k\ne j}|d_k|\le\epsilon. \tag{2}
\]
The signed-permutation assumptions make these inequalities sufficient as well as necessary: choose the matching permutation and signs, keeping the observed joint innovation-proxy law unchanged. This does not prove that the substantive bounded-exclusion assumption is true. It describes its implications within the declared statistical model.

For instance $d=(0.8,0.1,0)'$, $\epsilon=0.15$, $\kappa=0.5$ identifies the first positive column. The vector $(0.8,0.7,0)'$ with the same bounds makes every candidate fail, since either large coordinate would be a forbidden nontarget covariance. With $\epsilon=0.9$, $\kappa=0.5$, both first and second positive columns are admissible. These arithmetic examples illustrate rejection and ambiguity, not empirical proxy quality.

\section{Uncertainty without a discontinuous winner rule}
Suppose $B_0$ is known in a benchmark and an independent sample supplies $\widehat d$ with a simultaneous event
$\|\widehat d-d\|_\infty\le r$ of probability at least $1-\alpha$. An outer confidence set of labels retains $(j,s)$ when
\[
 s\widehat d_j+r\ge\kappa,\qquad
 \max_{k\ne j} (|\widehat d_k|-r)_+\le\epsilon. \tag{3}
\]
Every truly admissible label satisfies (3) on that event, so projecting the retained labels to their paths covers the entire population path set with probability at least $1-\alpha$. The separate inequalities may admit a candidate without a jointly feasible nuisance vector when other restrictions are added; this makes the set conservative, never anti-conservative. With a coordinate box and just (2), feasibility can be checked directly coordinate by coordinate.

When $B_0$ and the propagation matrices are estimated, their uncertainty must enter the same projection. A confidence region for $(B_0,A_0,\ldots,A_H,c)$ can be propagated through (1)--(3), excluding singular $B_0$ or allowing an uninformative result near singularity. Holding estimated impact directions fixed while varying only proxy correlations does not provide the claimed joint coverage. Covariance bounds based on concentration also require explicit moment or tail assumptions; arbitrary unbounded proxies do not satisfy a distribution-free Hoeffding bound.

A strict separation margin makes the labeling operation stable. If a true candidate satisfies its relevance inequality by more than $2r$ and all exclusion inequalities by more than $2r$, its empirical inequalities remain feasible. If every rival violates at least one inequality by more than $2r$, those rivals are rejected by (3). Without a positive margin, labels may remain multiple with arbitrarily large samples along local sequences. Reporting that ambiguity is the appropriate inferential behavior.

\section{Continuous ambiguity cannot be ignored}
Suppose two independent structural shocks are standard Gaussian and no higher-moment distinction separates them. Replacing their two columns by an arbitrary orthogonal rotation and rotating the shocks oppositely preserves their joint law. The compatible set then contains a circle, not merely signed permutations. Enumerating $2^m m!$ possibilities misses valid models. With $B_0=I_2$, the first coordinate's impact of a freely rotating unit column ranges over $[-1,1]$; a nonnegative-impact sign restriction reduces it to $[0,1]$, still a continuum.

A valid nonzero one-shock proxy can select a direction in that block, but its exclusion has the same economic burden as before. A continuous compact rotation set can be projected numerically using nets and continuity bounds, with relaxed restrictions to avoid missing feasible boundary points. Finite output is then an approximation to the set, not a proof of point identification. The same rotation must be retained across all horizons; choosing a favorable column separately at each horizon creates impossible response paths.

\section{Interpretation and unresolved content}
The finite algorithm completely resolves labeling conditional on a correct finite compatible set and declared restrictions. It does not identify which external instrument actually excludes all other economic shocks, or whether measured information omits anticipated policy. Statistical independence can distinguish latent directions while leaving their institutional interpretation open. A successful application needs external measurements and timing restrictions strong enough to separate economic stories, as well as a test of their compatibility. No policy label or causal response is inferred from statistical independence alone.

\section{Completion Estimate}
60\%

This is a post hoc, heuristic estimate for the full catalogue target.
The attempt gives exact signed-permutation label enumeration, proxy compatibility and confidence restrictions, and examples separating unique responses from unique economic labels. Full interpretation still requires handling continuous compatible rotations and validating external exclusions and institutional timing rather than assuming a finite statistical ambiguity set.

\begin{thebibliography}{9}
\bibitem{lewis} D. J. Lewis (2025), \emph{Identification Based on Higher Moments in Macroeconometrics}, Annual Review of Economics 17, 665--693, Section 6.4. \url{https://discovery.ucl.ac.uk/id/eprint/10208055/1/Lewis_annurev-economics-070124-051419.pdf}. The source discusses the shock-labeling issue; the finite procedures and illustrative moment vectors above state their own assumptions.
\end{thebibliography}

\end{document}
